\documentclass[10pt,a4paper]{amsart}
\usepackage[margin=19mm]{geometry}
\usepackage{amsmath,amssymb,mathtools}
\usepackage[scr=rsfs,cal=euler]{mathalfa}
\usepackage{xfrac}
\usepackage[unicode,hidelinks,bookmarks=false]{hyperref}
\newcommand{\ie}{i.e.\ }
\newcommand{\cf}{cf.\ }
\newcommand{\resp}{respectively\ }
\newcommand{\Subsection}{\subsection}
\providecommand{\nobreakdash}{\nobreak\hbox{-}}
\setlength{\emergencystretch}{2em}
\multlinegap0pt
\usepackage{bm}
\usepackage{bbm}
\usepackage{dsfont}
\usepackage{xspace}
\usepackage{cancel}
\usepackage{mathtools}
\usepackage{amsthm}
\makeatletter
\@ifundefined{theorem}{\newtheorem{theorem}{Theorem}}{}
\@ifundefined{lemma}{\newtheorem{lemma}[theorem]{Lemma}}{}
\makeatother
\newcommand\bbm{\begin{bmatrix}}
\newcommand\ebm{\end{bmatrix}}
\newcommand\la{\lambda}
\newcommand{\A}{\mathcal{A}}
\newcommand{\cA}{\mathcal{A}}
\newcommand{\inv}{^{-1}}
\title[Cauchy Transforms of Colored Graphs in Two Variables]{Cauchy Transforms of Colored Graphs in Two Variables}
\author{Lily Adlin, Giovani Thai, Samuel Tiscareno, Ryan Tully-Doyle}
\date{Source: arXiv:2410.10695v2 (CC BY 4.0)}
\begin{document}
\maketitle

\noindent\emph{Source and licence.} Cauchy Transforms of Colored Graphs in Two Variables. Version arXiv:2410.10695v2 (CC BY 4.0), distributed under the Creative Commons Attribution 4.0 International licence, as recorded in the arXiv licence metadata for this exact version and shown on the abstract page https://arxiv.org/abs/2410.10695v2. Full rights notice: licenses/arxiv-2410.10695-CC-BY-4.0.txt.\par\smallskip

\noindent\textit{Editorial scope.} Counted unit: the Theorem whose source label is \texttt{contactorder} in the article (source0.tex lines 1087--1089, source label \texttt{contactorder}), delivered together with the article's own complete original proof of it (source0.tex lines 1091--1167, one \texttt{proof} environment of 77 lines). No mathematics is written, repaired or re-derived: statement and proof are the article's own text line for line. Required context retained uncounted: relabel (lines 364--367); pathorder (lines 1076--1078). Every definition, display and label of those lines is retained unchanged. Omitted from this package and not redistributed: 1--363, 368--1075, 1079--1086, 1090--1090, 1168--1324. Nothing inside a retained excerpt is omitted or altered: every retained body is delivered exactly as the article's lines are written, so no omission receipt of any kind accompanies this package. Citations: the retained excerpts carry no citation command at all, so no bibliography entry of the article is transcribed and no reference is redistributed. Reference adaptation: none was needed and none was made; every reference inside the delivered unit resolves inside it, so no unresolved reference or citation is produced. Printed slips kept literal: no printed word, symbol, hypothesis or number of the counted statement or of its proof was altered. Layout and class: the article's own class is \texttt{amsart} with packages including amsmath, amsthm, amssymb, amsfonts, enumerate, nicematrix, enumitem, graphicx, hyperref, tikz; the wrapper is the lane's pinned \texttt{amsart} editorial wrapper at 10pt on A4, which restates the theorem-like environments the retained text needs and adds the article's own macro definitions that the retained text needs (\texttt{\textbackslash A}, \texttt{\textbackslash bbm}, \texttt{\textbackslash cA}, \texttt{\textbackslash ebm}, \texttt{\textbackslash inv}, \texttt{\textbackslash la}), with extra wrapper packages (bm, bbm, dsfont, xspace, cancel, mathtools). The delivered numbering is the unit's own. Assets: no retained span contains a figure environment, a graphics-inclusion command, a PDF-page inclusion, a table or a TikZ picture; no image file is copied into this package, the package declares no asset dependency and no image file is redistributed with it. Licence: version arXiv:2410.10695v2 of this article is distributed under the Creative Commons Attribution 4.0 International licence, as recorded in the arXiv licence metadata for this exact version and shown on the abstract page https://arxiv.org/abs/2410.10695v2. A full rights notice is delivered at licenses/arxiv-2410.10695-CC-BY-4.0.txt. Characters: every retained excerpt is pure ASCII, so the delivered engine, XeLaTeX with its default Latin Modern text font, reports no missing character.
\begin{lemma}\label{relabel}
    Let $G$ be a graph with colored adjacency matrix $\cA_G$ and $H$ be a graph with colored adjacency matrix $\mathcal{A}_H$. If $H = \phi(G)$ with graph isomorphism $\phi$ then 
    $$f^k_G(z,w) = f^{\phi(k)}_H(z,w).$$
\end{lemma}

\begin{lemma} \label{pathorder}
    In a graph $G$, if the shortest path from vertex $i$ to $j$ is of length $n$ then the order of vanishing of the path function at $\infty$ from $i$ to $j$ will be $n+1$.
\end{lemma}

\begin{theorem}\label{contactorder}
Let $G$ be a colored graph with one vertex, say $j$, labelled $w$ and all others $z$. Then the order of contact of the Nevanlinna representation at vertex $i$ is $2n$ where the minimum path length from $i$ to $j$ is $n$.
\end{theorem}

\begin{proof}
The argument goes by cases.

\textbf{Case 1:} First, consider the case where $n \geq 2$ (that is, the vertices of interest aren't directly connected). By Lemma \ref{relabel}, without loss of generality we consider the case where we calculate the Nevanlinna representation from vertex 1 and the $w$ is located at vertex 2. Then the matrix $\A_G$ can be written
\[
\A_G = \bbm \bbm -z & 0 \\ 0 & -w \ebm & A_{12} \\ A_{21} & A_{22} \ebm,
\]
where $A_{22}$ represents the graph formed from all paths from vertices adjacent to vertex $1$ to vertices adjacent to vertex $2$, and $A_{12} = A_{21}^T$ encodes the connection of the ``middle'' graph to the terminating vertices $1$ and $2$.
Taking a Schur complement with respect to the $(1,1)$ block gets
\[
f_G(z,w) = (\A_G\inv)_{(1,1)} = \bbm -z -g(z) & -p(z) \\ -p(z) & -w - h(z) \ebm\inv_{(1,1)}
\]
where $g$ is the sum of walk generating functions that represent walks leaving and returning to vertex 1 through each adjacent vertex, $h$ is the sum of walk generating functions representing walks leaving and returning to vertex $2$ through each vertex adjacent to vertex $2$, and $p$ is the sum of walk generating functions representing the paths connecting vertex $1$ to vertex $2$ through each pair of edges adjacent to vertex $1$ and vertex $2$ respectively. We have reduced the problem to finding the $(1,1)$ entry of the inverse of a $2 \times 2$ scalar matrix. Hence,
\begin{align*}
f_G(z,w) &= \bbm -z -g(z) & -p(z) \\ -p(z) & -w - h(z) \ebm\inv_{(1,1)} \\
&= \frac{-h-w}{-p^2+g h+g w+h z+w z}.
\end{align*}
To calculate the order of contact, we find the structure of the level curves in a neighborhood of $(\infty, 0)$ by setting $$f_G(z,w) = \la.$$ Solving for $w$ gives the level curve parametrization
\[
w = \Lambda_\la(z) = \frac{p^2 \la-g h \la-h \la z-h}{g \la+\la z+1} = -h + p^2 \cdot \frac{\la}{1 + \la(g + z)}.
\]


So the analysis comes down to understanding the function
\[
\frac{\la}{1 + \la(g(z) + z)}.
\]
As $g(z)$ is a sum of walk generating functions in $z$, $g$ has a power series expansion at $\infty$, and so expanding the rational expression as a series in $z$ at $\infty$, we get
\[
\frac{\la}{1 + \la(g(z) + z)}  = \frac{1}{z}  - \frac{a_0}{z^2} - \frac{1}{\la z^2}  + O\left(\left(\frac{1}{z}\right)^3\right).
\]
That is, the first term that depends on $\la$ is order 2. Note that $p$, another sum of walk generating functions, when expanded as a power series at $\infty$, has  lowest order term of order $n-1$ where $n$ is the shortest path length between $i$ and $j$ by Lemma \ref{pathorder}. So $p^2$ has lowest order term $2n - 2$. Thus 
\[
\Lambda_\la(z) = -h + p^2 \cdot \frac{\la}{1 + \la(g + z)} = -h + \frac{p^2}{z}  - \frac{a_0 p^2 }{z^2} - \frac{p^2}{\la z^2}  + O\left(\left(\frac{1}{z}\right)^3\right),
\]
which implies that generically $\Lambda_\la(z) - \Lambda_\mu(z)$ has order of vanishing $2n$ and hence $f_G$ has order of contact $2n$ at $(\infty, 0)$. 

\textbf{Case 2:} Now consider the case where the distance from vertex 1 to vertex 2 is 1 - that is, the $z$-vertex and the $w$-vertex of the Nevanlinna representation are connected. In this case, 
\[
\A_G = \bbm \bbm -z & 1 \\ 1 & -w \ebm & A_{12} \\ A_{21} & A_{22} \ebm.
\]
Then
\begin{align*}
f_G(z,w) = (\A_G)\inv_{(1,1)} &= \bbm -z -g(z) & 1-p(z) \\ 1-p(z) & -w - h(z) \ebm\inv_{(1,1)} \\
&= \frac{-h-w}{-p^2+2 p+g h+g w+h z+w z-1}.
\end{align*}
Choosing a value $\la$ and solving for $w$ as a function of $z$ for the level curve function gives
\[
w = \Lambda_\la(z) = -h + \frac{(p-1)^2 t}{g t+t z+1}
\]
Performing a similar expansion at $\infty$ gives that $\frac{1}{\la z^2}$ is the first term depending on $\la$, and so the order of contact is $2$.


\textbf{Case 3:} Finally, consider the case where we compute a Nevanlinna representation with respect to a vertex colored $w$, enumerated as vertex 1. Then we can write
    
\[
    \A_G = \begin{bmatrix}
        w & A_{12} \\
        A_{12}^T & A_{22}
    \end{bmatrix}.
\]
Now compute the Nevalinna representation to get

    \begin{align*}
        f_G(z,w) = (\A_G\inv)_{(1,1)} &= (-w - A_{12}A_{22}^{-1}A_{12}^T)^{-1} \\
        &= \dfrac{1}{-w-p(z)}
    \end{align*}

   where $p(z)$ is the sum of walk generating functions leaving and returning to vertex $1$ through adjacent vertices. We now compute $\Lambda_\la(z)$ as above, solving $f_G(z,w) = \la$ for $w$, which gives
   \[
   \Lambda_\la(z) = -p(z) - \frac{1}{\la}.
   \]

 By Lemma \ref{pathorder}, $p(z)$ has lowest order term at least 1, and so the constant term $-\dfrac{1}{\la}$ cannot be canceled out by $p(z)$. Then the order of contact is $0$. 
    

\end{proof}

\end{document}
